\section{近期研究（2026 年 8--11 月）}
\paragraph{Scene2Sound：将物体级音源绑定到可漫游的高斯场景}
2026.08~\cite{2608_00463}。Scene2Sound 面向已有的 3DGS 场景，先选取覆盖场景的观察位姿，使用视觉语言模型区分不需要定位的环境背景声与可定位的发声物体。模型将不同视图的物体掩码关联到高斯集合，再依据集合重合匹配三维实例，生成的音源交由对象式音频引擎按照听者位置渲染。

作者以 24 个生成的 SoundscapePLY 场景进行主要评测，并在由真实 360 度影像重建的 D-SAV360 场景检查迁移表现。19 位参与者以移动视点比较三种空间音频基线。论文的失败案例显示，如果分割阶段无法定位发声物体，系统可能只能合成无明确方位的背景声。

\paragraph{GS-Voxel：为大范围航拍场景构造结构化高斯潜变量}
2026.08~\cite{2608_17988}。GS-Voxel 将预优化且满足输入约定的高斯集合确定性组织为稀疏体素。因子化 VAE 分别编码体素占据和局部高斯属性，图像条件流模型在潜空间生成，带重叠的平铺推理将合成范围拓展到单块训练裁剪之外。这里避免的是逐场景向固定大小模板拟合，训练资料仍来自已重建的 3DGS 场景。

论文在各含 1.5 万张渲染图像的生成与留出参照集合上报告 FID 28.0、KID 0.020，并明确区分不同论文的相机协议和数据。当前实现只在 SH0 航拍场景开发和验证，改变卫星影像来源后的鲁棒性尚未评估。极密体素的容量限制和有限空间分辨率也可能分别造成原语丢弃与细结构退化。

\paragraph{SPAR3S：在稀疏体素潜空间自回归补全场景}
2026.09~\cite{2609_03931}。SPAR3S 将稀疏条件图像编码到与高斯解码器对齐的体素潜空间。掩码自回归模型在其中联合预测新增占据和潜变量，随后解码出可供新视点渲染的高斯集合。研究针对条件视图很少、彼此甚至没有重叠的室内场景重建。

论文在合成室内 3DFront 与真实视频 RealEstate10K 上评估，前者提供真值位姿并从渲染深度回投点图，后者由 MASt3R-SfM 估计初始几何。作者指出占据预测漏掉墙壁或物体体素时会留下空洞，有限体素分辨率和每体素的高斯密度也会引入局部模糊。

\paragraph{Filling the Unseen：分阶段视角外推与质量感知监督}
2026.09~\cite{2609_13262}。该方法先检查新相机视图的重叠和依赖，把可独立外推的相机视图分组并行处理，再补全依赖视图。质量感知掩码以视线方向相关的高斯渲染质量估计调节损失权重，使生成监督用于缺失或质量较低的区域，同时减少对可靠观测区域的干扰。

论文在 Mip-NeRF 360 的 4 个场景采用无真值代理指标，在 ScanNet++ 的 10 个场景将一半区域留作外推目标并以原影像评估。两套协议分别衡量感知代理和有图像参照的受控外推表现；在真实未知区域，代理指标不能证明生成内容与未观测事实吻合。

\paragraph{PhGS：对单视图前馈高斯先裁剪再细化}
2026.09~\cite{2609_20623}。PhGS 承接每像素预测多个原语的单视图前馈骨干。它先依重要性从预测的高斯集合裁剪，再用学习得到的循环模块直接更新保留原语，避免单纯删除原语带来的外观或几何损伤。混合不同保留比例训练让一个细化器可在推理时适配不同保留数目。

作者在 RealEstate10K 上训练和评测，并以 KITTI 和 DL3DV 做零样本跨域评测。在 Flash3D、Niagara、SHARP 三种骨干的域内实验中，主设置可舍弃一半预测高斯并保留有竞争力的渲染表现。细化模块仍需训练，跨域结论限于论文测试的骨干和数据。

\paragraph{GenNVS：把前景背景分别建模，再用三维先验约束视频外推}
2026.09~\cite{2609_34579}。单张图像外推新视图时，前景和背景不一定能由统一变形解释。GenNVS 先分别建立前景和背景，再对前景高斯作粗到细配准。从粗三维场景得到的重投影视图和有效掩码通过双流输入控制视频扩散，使它能区分有几何支撑的区域和需要补全的区域。

基线包含前馈与生成式新视图方法，原文在 DL3DV 和 Mip-NeRF 360 报告渲染与姿态一致性对照，并演示场景编辑。训练依赖配对视频资料和高算力，推理设置为 50 步去噪。补全区域的视觉稳定性仍不能作为场景真实结构的证据。

\paragraph{Edit4D：将静态背景与动态人物分开生成和编辑}
2026.10~\cite{2610_21668}。Edit4D 从文本分别生成高斯背景和动态人物，在统一场景中通过目标区域定位与接触约束完成组合。组件独立表示使替换人物外观或调整人物落点时可以保留已有背景。模型还可渲染改变相机位姿后的动态场景。

论文沿用 Free4D 的正面视频协议评估文本对齐、时序一致性和美学质量，并以人物外观编辑和场景重组对比必须全局重生成的 Free4D。作者记录了椅子被桌子遮挡时人物落点不准、与桌面几何相交的案例，因而放置的可靠性依赖目标可见程度。

\paragraph{TriFold：用规则三平面编码不规则高斯场景}
2026.10~\cite{2610_23050}。TriFold 将带有空间密度差异和可变原语数量的场景高斯集合编码为固定布局的三平面变分潜变量。编码采用单元查询及以原语为锚的跨平面注意力，解码器再按占据与局部原语置信预测分配高斯。规则潜空间也供条件流匹配生成使用。

论文在 InteriorGS 与 SceneSplat-49K 的 Aria Synthetic Environments 子集验证重建，按原场景划分训练与测试以避免裁剪泄漏。论文另对点云条件和单图全形状补全展示生成结果，并指出观察稀疏的场景边界仍可能较真值模糊。

\paragraph{Plant Orbiter：用相机轨迹已知的生成视频构造植物高斯}
2026.11~\cite{2611_12739}。Plant Orbiter 由单张植物图像合成 73 帧、围绕静止植物的短视频，借预定轨道估计相机位姿并拟合可对齐输入的高斯模型。前接文生图模型后可用同一流程生成文本指定的植物资产。

论文报告各测试集的最佳几何 F-score，并在由两株植物、17 个时间点组成的 GrowSplat 样本上用点集测量植株高度，结论报告平均高度误差 5.9\%。单视图看不到的背侧和遮挡部位由模型推断，固定高度单轨环绕也会造成顶视或底视伪影，限制其用于强自遮挡结构的表型测量。

\paragraph{VideoRNG：用标准化视频观测生成可重照明的神经高斯}
2026.11~\cite{2611_25338}。VideoRNG 在规定的转台相机与照明下从文本或图像合成观察序列，将 29 帧编码至 8 个视频潜帧，前景特征经重建网络一次预测高斯几何与神经外观。几何与外观的组织允许在新观察方向和新光照下渲染资产。

作者用约 2.2 万个筛选后的 TexVerse 纹理资产训练，并在留出的视点和灯光条件下评估重建。论文深度锚定对比的 PSNR 为 29.22 dB，联合深度监督为 25.39 dB，其含义限于该比较。生成视频的跨帧不一致可在重建中留下重影，尖锐镜面效果仍难复现。
